\section{Cinematica de la particula}

La cinemática es la rama de la mecánica que estudia el movimiento de las partículas sin considerar las causas que lo originan. Para describir dicho movimiento, se introducen magnitudes fundamentales como la posición, la velocidad y la aceleración, todas ellas dependientes del tiempo.
\vspace{-20pt}

\begin{align*}
\vec{v_m} &= \lim_{\Delta t \to 0} \left(\frac{\Delta \vec{r}}{\Delta t}\right)\\
\Aboxed{\vec{v_m} &= \frac{d\vec{r}}{dt}(t)}\\[10pt]
\vec{a_m} &= \lim_{\Delta t \to 0} \left(\frac{\Delta \vec{v}}{\Delta t}\right)\\
\Aboxed{\vec{a} &= \frac{d\vec{v}}{dt}(t)}
\end{align*}

Las relaciones anteriores también pueden expresarse en forma integral, lo que permite determinar la velocidad y la posición a partir de la aceleración:

\vspace{-20pt}

\begin{align*}
v &= \frac{dr}{dt}(t) \\
\int dr &= \int v(t)\, dt \\
a &= \frac{dv}{dt}(t) \\
\int dv &= \int a(t)\, dt
\end{align*}
